\section{PolyArt Dataset}

\begin{figure*}[ht]
    \centering\includegraphics[width=\linewidth]{images/polyart_pipeline.jpg}
    \vspace{-5mm}
    \caption{Overview of the PolyArt dataset construction pipeline. Each painting is processed through four specialized vision-language models to produce independently generated perspectives (narrative, formal, emotional, historical), which are then harmonized into a unified caption.}
\label{fig:polyart_framework}
\end{figure*}

We introduce PolyArt, a large-scale multi-perspective dataset grounding visual art understanding in four distinct interpretive dimensions.
This section describes the data foundation, the four-perspective architecture, the model selection and generation pipeline, quality control procedures, and the two benchmark tasks built on the dataset.

\subsection{Data Foundation}

PolyArt is built on the WikiArt dataset~\cite{2017omniart}, which covers approximately 75,921 artworks across 27 style categories spanning the 15th through the 21st century.
After filtering to artworks with verified image files on disk, the dataset comprises approximately 70,000 paintings from over 1,100 artists.
Each painting is associated with structured metadata including title, artist name, production date, style classification, and artist school.

For a subset of artworks, we additionally retrieve artist-level biographical and contextual information from Wikidata, including movement affiliations, key biographical events, and cross-references to other artists.
This Wikidata context is used as a retrieval source for grounding the historical and contextual perspective.

\subsection{Four-Perspective Architecture}

Expert art interpretation operates simultaneously across multiple analytical registers.
Rather than generating a single descriptive caption, PolyArt decomposes each artwork's annotation into four independently generated perspectives:

\begin{equation}
P(w) = \{ e^{\text{narr}},\ e^{\text{form}},\ e^{\text{emot}},\ e^{\text{hist}} \},
\end{equation}

where $w$ denotes a painting and each component captures a distinct interpretive dimension.

\paragraph{Narrative and Scene Interpretation ($e^{\text{narr}}$).}
A factual account of the depicted entities, figures, scene composition, and visual elements as they appear in the image.
This perspective answers \textit{what} is shown, without invoking symbolic or historical interpretation.

\paragraph{Formal Visual Analysis ($e^{\text{form}}$).}
An analysis of the painting's compositional structure, including spatial organization, color palette, brushwork, use of light and shadow, and visual rhythm.
This perspective corresponds to the vocabulary of formal art criticism, focusing on \textit{how} the work is constructed.

\paragraph{Emotional Response ($e^{\text{emot}}$).}
An affective characterization of the painting's perceived mood, atmosphere, and psychological tone as registered by a viewer.
This perspective captures the phenomenological dimension of art encounter---\textit{what it feels like} to look at the work.

\paragraph{Historical and Contextual Analysis ($e^{\text{hist}}$).}
An art-historical interpretation situating the work within its cultural context: movement affiliation, iconographic codes, period-specific symbolic meaning, and relevant biographical or historical circumstances.
This perspective answers \textit{why} the work looks and means as it does.

In addition to the four independent perspectives, each painting receives a \textit{unified caption} $e^{\text{unif}}$ produced by harmonizing all four into a single coherent description.
This unified caption serves as a strong single-text baseline in benchmark evaluations.

\subsection{Generation Pipeline}

A central design principle of PolyArt is the use of specialized models per perspective, each chosen for its demonstrated strength in the corresponding interpretive task.
Generating all perspectives with a single general-purpose model risks homogenizing the output; using specialized models encourages each perspective to reflect genuine domain fit.

\paragraph{Narrative perspective.}
Narrative captions are generated using \textbf{Qwen2.5-VL-72B-Instruct}~\cite{qwen25vl}, a large-scale vision-language model with strong scene understanding.
The 72B scale supports detailed factual description while the model's training on diverse visual corpora reduces systematic bias toward art-specific jargon, keeping the output grounded in observable visual content.

\paragraph{Formal analysis.}
Formal perspectives are generated using \textbf{GalleryGPT}~\cite{MM24GalleryGPT}, a LLaVA-7B model fine-tuned on the PaintingForm dataset specifically for formal art analysis.
This specialization ensures that output adheres to the vocabulary and analytical conventions of formal criticism rather than defaulting to general image description.

\paragraph{Emotional response.}
Emotional perspectives are generated using \textbf{Qwen3-VL-8B-Instruct}, conditioned on the painting image together with human-written emotional utterances from the ArtEmis dataset~\cite{achlioptas2021artemis} as few-shot context.
ArtEmis provides on average 5.7 affective utterances per WikiArt painting, grounding the emotional generation in authentic human responses rather than model inference alone.

\paragraph{Historical and contextual analysis.}
Historical perspectives are generated using \textbf{Qwen2.5-VL-72B-Instruct} augmented with a Wikidata knowledge graph retrieved via LightRAG~\cite{lightrag2024} in hybrid retrieval mode.
For each artwork, relevant Wikidata snippets covering the artist biography, movement, and historical events are retrieved and injected into the generation context.
This retrieval-augmented approach grounds historical claims in verifiable external knowledge and reduces hallucination of biographical or period-specific facts.

\paragraph{Unified caption.}
Given the four independently generated perspectives $\{e^{\text{narr}}, e^{\text{form}}, e^{\text{emot}}, e^{\text{hist}}\}$, a unified caption $e^{\text{unif}}$ is produced via \textbf{Qwen3-8B}~\cite{qwen3} using a harmonization prompt that synthesizes the four perspectives into a single coherent, non-redundant description of approximately 150 words.
The unified caption preserves interpretive breadth while eliminating cross-perspective repetition, and is designed for downstream tasks such as retrieval and language-model grounding where a holistic single-text representation is preferred.

Formally, the generation of perspective $v$ for painting $w$ with image $I$, metadata $M$, and optional retrieved context $\mathcal{S}$ is:

\begin{equation}
e^{v} = \text{VLM}_v\!\left(I,\ M,\ \mathcal{S},\ \pi_v\right),
\end{equation}

where $\pi_v$ is a perspective-specific prompt template designed to enforce interpretive focus and minimize cross-perspective overlap.

\subsection{Quality Control}

\paragraph{Cross-view consistency.}
We enforce consistency between the historical perspective and the WikiArt style classification label.
Let $m_w$ denote the known style label for painting $w$.
We require semantic alignment between $e^{\text{hist}}$ and $m_w$:

\begin{equation}
\text{sim}\!\left(e^{\text{hist}},\ m_w\right) \ge \tau,
\end{equation}

where $\text{sim}(\cdot, \cdot)$ is sentence-level cosine similarity computed via a pretrained text encoder and $\tau$ is a threshold calibrated on the gold subset.
Perspectives failing this check are flagged for regeneration under constrained prompting.

\paragraph{Gold subset evaluation.}
We evaluate perspective quality via two complementary checks reported in Section~\ref{sec:quality}.
First, a \textit{metadata factual verification} checks whether $e^{\text{hist}}$ correctly identifies the artist name, art movement, and production period against WikiArt metadata across all 74,234 paintings.
Second, an \textit{LLM-as-judge} evaluation uses Gemma-3-27B~\cite{gemma3}---a vision-language model from a different model family than all generation models---to rate all five perspectives on perspective fidelity, factual accuracy, and depth on a stratified 300-painting sample.
Using an independent judge from a different model family mitigates self-evaluation bias.

\paragraph{Dataset-level statistics.}
Table~\ref{tab:dataset_stats} summarizes the final dataset statistics, including coverage per perspective and comparison against prior art captioning datasets.

\begin{table}[h]
\centering
\caption{PolyArt dataset statistics and comparison with related datasets.}
\label{tab:dataset_stats}
\begin{tabular}{lrrl}
\toprule
\textbf{Dataset} & \textbf{Paintings} & \textbf{Captions/img} & \textbf{Perspectives} \\
\midrule
SemArt~\cite{semart_2018}        & 21,384 & 1 & Historical/iconographic \\
ArtEmis~\cite{achlioptas2021artemis}       & 81,446 & 5.7 & Emotional \\
ExpArt~\cite{ExpArt_acl2024}          & 19,181 & 1 & Formal/historical \\
GalleryGPT~\cite{MM24GalleryGPT}     & 75,336 & 1 & Formal \\
\midrule
\textbf{PolyArt (ours)} & \textbf{$\sim$70,000} & \textbf{5} & \textbf{All four + unified} \\
\bottomrule
\end{tabular}
\end{table}

\subsection{Benchmark Task: Perspective Complementarity}

PolyArt includes a benchmark for evaluating whether the four perspectives encode genuinely distinct visual information.

\textbf{Question}: Do the four perspectives encode genuinely distinct information, or are they stylistic paraphrases of the same description?

\textbf{Setup}: Given text from one or more perspectives (without the original image), a text-to-image model generates a reconstruction of the original painting.
We evaluate ten perspective conditions: four singles ($e^{\text{narr}}$, $e^{\text{form}}$, $e^{\text{emot}}$, $e^{\text{hist}}$), four leave-one-out triples (each omitting one perspective), the full four-perspective combination, and the unified caption as a holistic baseline.
Each reconstruction is compared to the original across three fidelity dimensions:

\begin{equation}
\mathcal{F}(w, \hat{w}) = \bigl(\delta_{\text{style}},\ \delta_{\text{comp}},\ \delta_{\text{emot}}\bigr),
\end{equation}

where $\delta_{\text{style}}$ is CLIP ViT-L/14 cosine similarity, $\delta_{\text{comp}}$ is DINOv3-Large CLS token cosine similarity capturing spatial structure and composition, and $\delta_{\text{emot}}$ is CLIP zero-shot emotion agreement against nine ArtEmis emotion labels.

\textbf{Hypothesis}: Each perspective recovers a different subset of fidelity dimensions, and reconstruction improves as perspectives are combined.
If perspectives were redundant, adding more would not improve reconstruction fidelity.

\textbf{Scale}: Evaluated on a stratified sample of 1,350 paintings (50 per style $\times$ 27 styles) using two text-to-image models---FLUX.2-Klein-4B~\cite{flux2} and Qwen-Image-2512~\cite{qwenimage}---to confirm that findings are model-agnostic.

